% Options for packages loaded elsewhere
% Options for packages loaded elsewhere
\PassOptionsToPackage{unicode}{hyperref}
\PassOptionsToPackage{hyphens}{url}
%
\documentclass[
  ignorenonframetext,
]{beamer}
\newif\ifbibliography
\usepackage{pgfpages}
\setbeamertemplate{caption}[numbered]
\setbeamertemplate{caption label separator}{: }
\setbeamercolor{caption name}{fg=normal text.fg}
\beamertemplatenavigationsymbolsempty
% remove section numbering
\setbeamertemplate{part page}{
  \centering
  \begin{beamercolorbox}[sep=16pt,center]{part title}
    \usebeamerfont{part title}\insertpart\par
  \end{beamercolorbox}
}
\setbeamertemplate{section page}{
  \centering
  \begin{beamercolorbox}[sep=12pt,center]{section title}
    \usebeamerfont{section title}\insertsection\par
  \end{beamercolorbox}
}
\setbeamertemplate{subsection page}{
  \centering
  \begin{beamercolorbox}[sep=8pt,center]{subsection title}
    \usebeamerfont{subsection title}\insertsubsection\par
  \end{beamercolorbox}
}
% Prevent slide breaks in the middle of a paragraph
\widowpenalties 1 10000
\raggedbottom
\AtBeginPart{
  \frame{\partpage}
}
\AtBeginSection{
  \ifbibliography
  \else
    \frame{\sectionpage}
  \fi
}
\AtBeginSubsection{
  \frame{\subsectionpage}
}
\usepackage{iftex}
\ifPDFTeX
  \usepackage[T1]{fontenc}
  \usepackage[utf8]{inputenc}
  \usepackage{textcomp} % provide euro and other symbols
\else % if luatex or xetex
  \usepackage{unicode-math} % this also loads fontspec
  \defaultfontfeatures{Scale=MatchLowercase}
  \defaultfontfeatures[\rmfamily]{Ligatures=TeX,Scale=1}
\fi
\usepackage{lmodern}

\usetheme[]{metropolis}
\ifPDFTeX\else
  % xetex/luatex font selection
\fi
% Use upquote if available, for straight quotes in verbatim environments
\IfFileExists{upquote.sty}{\usepackage{upquote}}{}
\IfFileExists{microtype.sty}{% use microtype if available
  \usepackage[]{microtype}
  \UseMicrotypeSet[protrusion]{basicmath} % disable protrusion for tt fonts
}{}
\makeatletter
\@ifundefined{KOMAClassName}{% if non-KOMA class
  \IfFileExists{parskip.sty}{%
    \usepackage{parskip}
  }{% else
    \setlength{\parindent}{0pt}
    \setlength{\parskip}{6pt plus 2pt minus 1pt}}
}{% if KOMA class
  \KOMAoptions{parskip=half}}
\makeatother


\usepackage{longtable,booktabs,array}
\usepackage{calc} % for calculating minipage widths
\usepackage{caption}
% Make caption package work with longtable
\makeatletter
\def\fnum@table{\tablename~\thetable}
\makeatother
\usepackage{graphicx}
\makeatletter
\newsavebox\pandoc@box
\newcommand*\pandocbounded[1]{% scales image to fit in text height/width
  \sbox\pandoc@box{#1}%
  \Gscale@div\@tempa{\textheight}{\dimexpr\ht\pandoc@box+\dp\pandoc@box\relax}%
  \Gscale@div\@tempb{\linewidth}{\wd\pandoc@box}%
  \ifdim\@tempb\p@<\@tempa\p@\let\@tempa\@tempb\fi% select the smaller of both
  \ifdim\@tempa\p@<\p@\scalebox{\@tempa}{\usebox\pandoc@box}%
  \else\usebox{\pandoc@box}%
  \fi%
}
% Set default figure placement to htbp
\def\fps@figure{htbp}
\makeatother





\setlength{\emergencystretch}{3em} % prevent overfull lines

\providecommand{\tightlist}{%
  \setlength{\itemsep}{0pt}\setlength{\parskip}{0pt}}



 


\setbeamerfont{title}{size=\large} \usepackage{hyperref} \setbeamertemplate{footline}[frame number]
\makeatletter
\@ifpackageloaded{caption}{}{\usepackage{caption}}
\AtBeginDocument{%
\ifdefined\contentsname
  \renewcommand*\contentsname{Table of contents}
\else
  \newcommand\contentsname{Table of contents}
\fi
\ifdefined\listfigurename
  \renewcommand*\listfigurename{List of Figures}
\else
  \newcommand\listfigurename{List of Figures}
\fi
\ifdefined\listtablename
  \renewcommand*\listtablename{List of Tables}
\else
  \newcommand\listtablename{List of Tables}
\fi
\ifdefined\figurename
  \renewcommand*\figurename{Figure}
\else
  \newcommand\figurename{Figure}
\fi
\ifdefined\tablename
  \renewcommand*\tablename{Table}
\else
  \newcommand\tablename{Table}
\fi
}
\@ifpackageloaded{float}{}{\usepackage{float}}
\floatstyle{ruled}
\@ifundefined{c@chapter}{\newfloat{codelisting}{h}{lop}}{\newfloat{codelisting}{h}{lop}[chapter]}
\floatname{codelisting}{Listing}
\newcommand*\listoflistings{\listof{codelisting}{List of Listings}}
\makeatother
\makeatletter
\makeatother
\makeatletter
\@ifpackageloaded{caption}{}{\usepackage{caption}}
\@ifpackageloaded{subcaption}{}{\usepackage{subcaption}}
\makeatother

\usepackage{bookmark}
\IfFileExists{xurl.sty}{\usepackage{xurl}}{} % add URL line breaks if available
\urlstyle{same}
\hypersetup{
  pdfauthor={Giorgio Arcara},
  hidelinks,
  pdfcreator={LaTeX via pandoc}}


\author{Giorgio Arcara}
\date{}

\begin{document}


\begin{frame}
% TITLE SLIDES
\title{Metodologia della Ricerca Clinica\\ \vspace{1em} \emph{10 - Interpretazione dei test}}
\author{Giorgio Arcara,\\ Università di Padova \\ IRCCS San Camillo, Venezia}

\titlegraphic{

\vspace*{7cm}
\includegraphics[scale=0.15]{Figures/LogoSanCamilloIRCCS_Unipd_alpha.png}
\begin{center}
\vspace{-2.2em}
\includegraphics[scale=0.15]{Figures/CC_license_3_0.png}
\end{center}
}
%\date{\today}
\maketitle
\end{frame}

\section{Interpretazione dei test}\label{interpretazione-dei-test}

\begin{frame}{Premessa}
\phantomsection\label{premessa}
\small

Quando si studiano gli aspetti statistici dei test in neuropsicologia
forense, emerge una criticità.

Molto spesso aspetti statistici (es. Validità, Affidabilità, ma anche
Errore di 1° Tipo, 2° Tipo), ci dicono informazioni sui test in generale
oppure su che percentuali di errore ci aspetteremmo se facessimo
infinite ripetizioni e con specifiche assunzioni.

Nella pratica di neurpsicologia clinica e forense, l'interesse è però
nel dare una risposta specifica, spesso puntuale relativamente ad uno
specifico individuo, con il massimo di accuratezza o precisione
possibile: ha un deficit cognitivo? Sta simulando un deficit?

L'aspetto chiave che tratteremo è che rispondere a queste domande non
può prescindere da un'interpretazione del dato, che tenga conto non solo
delle evidenze del test, ma dalle altre evidenze a disposizione.
\end{frame}

\begin{frame}{Interpretazione}
\phantomsection\label{interpretazione}
\begin{center}
\begin{tikzpicture}
    % Include your image
    \node[anchor=south west, inner sep=0] (image) at (0,0)
        {\includegraphics[width=0.8\textwidth]{Figures/Test_Methodology_transp.png}};

    % Set coordinate system to match the image
    \begin{scope}[x={(image.south east)}, y={(image.north west)}]
        
        % Draw a red rectangle (thin border)
        %\draw[red, thick, rounded corners] (0.25,0.3) rectangle (0.45,0.5);
        % Optional label
        %\node[red, font=\bfseries, above] at (0.35,0.5) {Region 1};
        
        %\pause
        % a red rectangle (thicker border)
        \draw[red, ultra thick] (0.2, 0.01) rectangle (0.75,0.32);


        
    \end{scope}
\end{tikzpicture}
\end{center}
\end{frame}

\begin{frame}{Interpretazione non vuol dire che va bene tutto}
\phantomsection\label{interpretazione-non-vuol-dire-che-va-bene-tutto}
Si potrebbe pensare che un accento sull' ``interpretazione'' vuol dire
un accento sulla soggettività del valutatore.

\begin{center}
\Large
\pause
\textbf{SBAGLIATO}
\end{center}
\pause
\normalsize
Per interpretazione si intende \underline{integrazione di tutte le evidenze disponibili} per raggiungere una conclusione sullo stato cognitivo del paziente.
Quello che si intende è dunque che non si può usare una serie di sillogismi per desumere deficit (se punteggio sotto cut-off → deficit), ma che anche il concludere un deficit è un interpretazione delle evidenze. I punteggi ai test rappresentano parte delle evidenze.
\end{frame}

\begin{frame}{Interpretazioni}
\phantomsection\label{interpretazioni}
Alcune interpretazioni sono migliori di altre.

Le interpretazioni più adeguate sono quelle che integrano tutte le
informazioni disponibili.
\end{frame}

\begin{frame}{Perché «interpretazione» e non ``lettura'' dei risultati?}
\phantomsection\label{perchuxe9-interpretazione-e-non-lettura-dei-risultati}
\small

Se un test possiede dati a sostegno di validità che misuri un certo
costrutto, possiamo ragionevolmente credere che tramite il test
misureremo quel costrutto e dovremmo sempre partire con questo
presupposto. Al contempo dobbiamo essere consapevoli che ci sono varie
occasioni in cui il numero non riflette ciò che dovrebbe riflettere.
Inoltre è noto che ogni comportamento è il frutto dell'interazione di
più funzioni cognitive (Lezak, 2000)

Ci sono però numerosi casi in cui il test potrebbe perdere la validità e
misurare altro.

\begin{itemize}
\tightlist
\item
  Il contesto della misurazione neuropsicologica è particolare proprio
  perché indaga un sistema che è potenzialmente danneggiato in maniera
  complessa.\\
\item
  nei casi di valutazione forense si aggiunge la complessità della
  possibile simulazione di deficit cognitivi.
\end{itemize}
\end{frame}

\begin{frame}{Un esempio di interpretazione (1/2)}
\phantomsection\label{un-esempio-di-interpretazione-12}
\small

\emph{Supponiamo di avere un paziente, GH, che effettua al test tra cui
il test di Linguaggio Figurato, in cui si indaga la capacità di
comprendere utilizzo di linguaggio in metafore, idiomi e proverbi (da
APACS, Arcara e Bambini, 2016). GH va effettivamente al sotto del
cut-off di normalità (10 su 15).}

\emph{Scenario 1: studiando la storia personale di GH, emerge che GH è
noto tra gli amici per il suo spiccato senso di umorismo, tanto che
spesso al bar si ferma per raccontare anche barzellette agli amici.
Varie persone che lo conoscono confermano questa voce sostenendo che
nonostante non fosse molto colto, a GH piaceva molto prendere gli altri
in giro usando giochi di parole.}
\end{frame}

\begin{frame}{Un esempio di interpretazione (2/2)}
\phantomsection\label{un-esempio-di-interpretazione-22}
\small

\emph{GH (si veda esempio), effettua una serie di test su comprensione
di linguaggio figurato in cui va effettivamente al sotto della norma (10
su 15).}

\emph{Scenario 2: studiando la storia personale di GH emerge che i suoi
problemi di comprensione del linguaggio erano molto comuni e noti. Al
bar spesso gli amici lo prendevano in giro proprio perché sembrava
spesso non capire cosa gli dicevano e spesso per questa ragione perdeva
la pazienza.}
\end{frame}

\begin{frame}{Esempi di interpretazione}
\phantomsection\label{esempi-di-interpretazione}
\small

L'accento sull'interpretazione indica che bisogna considerare le
evidenze a disposizione al di là del punteggio al test.

Nello scenario 2 difficilmente un'interpretazione sensata potrà
suggerire che ci sono evidenze di simulazione. Lo scenario 1 invece
sembra in qualche modo suggerirlo. È compito del neuropsicologo
raccogliere in maniera attiva le evidenze rilevanti e raggiungere una
conclusione.

Spesso (se non sempre) conclusione comporterà l'aggiungere qualcosa, non
il dedurre, che rappresenta appunto l' \emph{interpretazione} fatta
dalla/dal neuropsicologa/o.
\end{frame}

\begin{frame}{Non tutte le evidenze sono uguali}
\phantomsection\label{non-tutte-le-evidenze-sono-uguali}
\begin{center}
\includegraphics[width=0.15\linewidth,height=\textheight,keepaspectratio]{Figures/triangle.png}
\end{center}

Dire che è importante raccogliere evidenze non vuol dire che tutte le
evidenze hanno lo stesso peso e sono uguali.

\textbf{Saranno sempre più forti (e convincenti) quelle evidenze basate
su dati oggettivi}, su test, su misurazioni, rispetto ad aspetti
qualitativi e opinioni personali (es. ``secondo me sta comunque
fingendo'')
\end{frame}

\begin{frame}{Non tutte le evidenze sono uguali}
\phantomsection\label{non-tutte-le-evidenze-sono-uguali-1}
\begin{center}
\includegraphics[width=0.15\linewidth,height=\textheight,keepaspectratio]{Figures/triangle.png}
\end{center}

Per ogni insieme di evidenze sono possibili più interpretazioni.

Non esiste una formula matematica per dire come combinarle, ma esistono
interpretazioni peggiori ed interpretazioni migliori che dipendono da
come sono correttamente considerate le evidenze disponibili.

In molti casi (specie forensI) non sarà facile dire che si ha
un'interpretazione ``giusta'' in mano, ma solo un'interpretazione
supportata dalle evidenze.
\end{frame}

\begin{frame}{Alcune interpretazioni}
\phantomsection\label{alcune-interpretazioni}
Nelle slides che seguono saranno portati alcuni esempi di
interpretazioni di aspetti che sono molto comuni.

\begin{itemize}
\tightlist
\item
  Interpretazione di un risultato sotto cut-off come deficit cognitivo
\item
  Interpretare pattern di risultati di diversi test al di sotto della
  norma
\item
  Integrare le informazioni per interpretare i test
\end{itemize}
\end{frame}

\begin{frame}{\normalsize Interpretazione di un risultato sotto cut-off
come deficit cognitivo}
\phantomsection\label{interpretazione-di-un-risultato-sotto-cut-off-come-deficit-cognitivo}
\small

Se il mio soggetto va sotto cut-off i in un test di memoria a breve
termine e concludo che è presente di un danno cognitivo non è semplice
«lettura» dei risultati, ma è già \textbf{diagnosi descrittiva}.

Il neuropsicologo sta assumendo:

\begin{itemize}
\tightlist
\item
  Che il test abbia misurato ciò che intendeva misurare \pause
\item
  Che l'ansia abbia avuto un ruolo trascurabile \pause
\item
  Che il soggetto era sufficientemente motivato e non ha finto un
  deficit. \pause
\item
  Che eventuali danni ad altre funzioni cognitive non hanno avuto un
  ruolo cruciale. \pause
\item
  Che i dati normativi erano adeguati per stimare la prestazione del
  soggetto. \pause
\end{itemize}

Molte di queste assunzioni sono ``implicite'' e un buon neuropsicologo
(si veda principio 5 Interpretative Approach, in 1.Introduzione), ne è
consapevole
\end{frame}

\begin{frame}{\normalsize Interpretazione di un risultato sotto cut-off
come deficit cognitivo}
\phantomsection\label{interpretazione-di-un-risultato-sotto-cut-off-come-deficit-cognitivo-1}
\small

È importante sottolineare che il fare attenzioni alle assunzioni
implicite non significa dubitare di tutto.

Significa non sottovalutare quelle assunzioni che sono particolarmente
critiche perché potrebbero non essere valide.

Esempi: \pause

\begin{itemize}
\tightlist
\item
  in contesto forense non possiamo dare per scontato che il paziente non
  stia simulando.\pause
\item
  in vari casi non possiamo dare per scontato che i dati normativi siano
  adeguati.\pause
\item
  non possiamo dare per scontato che altre funzioni cogntive non di
  interesse non abbiamo influenzato la performance in un paziente con
  sospetto danno cerebrale.
\end{itemize}
\end{frame}

\begin{frame}{\normalsize Interpretazione di un risultato sotto cut-off
come deficit cognitivo}
\phantomsection\label{interpretazione-di-un-risultato-sotto-cut-off-come-deficit-cognitivo-2}
L'interpretazione di un test dorebbe essere fatta al termine della
valutazione, quando tutte le altre informazioni (anche provenienti da
altri test) sono disponibili

Uno degli errori più comuni è quello di interpretare già come
\emph{definitiva}, l'evidenza di un deficit già al primo risultato
ottenuto. Questo non vale solo per deficit, ma per qualsiasi altro tipo
di informazione raccolta.
\end{frame}

\begin{frame}{\normalsize Interpretare pattern di risultati di diversi
test al di sotto della norma}
\phantomsection\label{interpretare-pattern-di-risultati-di-diversi-test-al-di-sotto-della-norma}
Cosa concludere se più test sono sotto il cut-off di normalità?

Più test deficitari possono essere legati ad un singolo deficit
cognitivo (es. attenzione) o a deficit multipli. \pause

È difficile trovare dei principi logici per rispondere a come
interpretare queste situazioni. \pause

La/Il neuropsicologa/o deve integrare le informazioni a sua disposizione
in maniera sensata.
\end{frame}

\begin{frame}{\normalsize Integrare l'interpretazione di risultati a
test su costrutti a diversi livelli}
\phantomsection\label{integrare-linterpretazione-di-risultati-a-test-su-costrutti-a-diversi-livelli}
Attenzione all'integrazione di costrutti su diversi livelli. Devono
essere fatte in maniera sensata, tenendo conto della diversa gerarchia
dei costrutti misurati.

Può non avere senso concludere che il soggetto abbia un deficit in un
costrutto più generico e si hanno altre informazioni più dettagliate.

\pause

\emph{Esempio: in APACS, il paziente ha un deficit nel punteggio di Comprensione Pragmatica, ma ha fatto errori solamente nel test di Umorismo, non ha senso concludere generalmente che c’è un deficit di Comprensione. Abbiamo già elementi per un’interpretazione più specifica.}
\end{frame}

\begin{frame}{Integrare le informazioni per interpretare i test}
\phantomsection\label{integrare-le-informazioni-per-interpretare-i-test}
\begin{center}
\includegraphics[width=0.15\linewidth,height=\textheight,keepaspectratio]{Figures/triangle.png}
\end{center}

I punteggi ai test, il risultato del confronto con dati normativi e
cut-off sono \textbf{informazioni} che vanno interpretate con le altre
informazioni disponibili.
\end{frame}

\begin{frame}{Integrare le informazioni per interpretare i test}
\phantomsection\label{integrare-le-informazioni-per-interpretare-i-test-1}
\begin{itemize}
\tightlist
\item
  Risultati ai test
\item
  Informazioni dall'Intervista preliminare
\item
  Dati qualitativi emersi durante la valutazione
\item
  Informazioni dai familiari
\end{itemize}

Sono tutti elementi da integrare per raggiungere una \textbf{diagnosi
descrittiva} (quali deficit) sullo stato cogntivo del paziente e la
\textbf{diagnosi di compatibilità con un'eziologia} (quale eziologia)
dell'eventuale disturbo neuropsicologico.

In ambito forense occorre inoltre fare ipotesi su presunta
\textbf{causalità} rispetto all'evento rilevante.
\end{frame}

\begin{frame}{Informazioni nella valutazione neuropsicologica}
\phantomsection\label{informazioni-nella-valutazione-neuropsicologica}
Considerando tutte le evidenze raccolte nella valutazioni informazioni
che possono aiutare nella diagnosi sia i test, sia l' \emph{osservazione
clinica} hanno grande rilevanza e devono essere integrate per
raggiungere conclusioni sensate. Aspetti psicometrici dei test (loro
affidabilità, validità, etc.)
\end{frame}

\begin{frame}{Interpretare casi ambigui}
\phantomsection\label{interpretare-casi-ambigui}
Caso 1: Prestazioni ai limiti della norma. \pause \emph{Come
interpretare questi risultati?} \pause

Dipende. Non esiste una regola per interpretare i risultati ai limiti
della norma. L'interpretazione dipende dall'integrazione con tutte le
altre informazioni disponbili. Ricordarsi che i cut-off ci servono per
stimare la prestazione del soggetto
\end{frame}

\begin{frame}{Interpretare casi ambigui}
\phantomsection\label{interpretare-casi-ambigui-1}
Caso 2: Soggetto che ha età o scolarità al limite tra due fasce.

\emph{Come interpretare questi risultati, con che fascia confrontare?}

Il confronto con i dati normativi ci serve per stimare la prestazione
del soggetto. Se il soggetto ha età o scolarità (o altre variabili) a
cavallo tra due fasce probabilmente la cosa migliore è confrontare la
prestazione con entrambi i cut-off. La/Il neuropsicologa/o dovrà
integrare questare informazione con le altre disponibili per giungere ad
una conclusione.
\end{frame}

\begin{frame}{Un esempio da contesto forense (1/2)}
\phantomsection\label{un-esempio-da-contesto-forense-12}
\scriptsize

Siete consulenti tecnici per il caso di un individuo che lamenta un
deficit cognitivo in seguito ad un incidente sul lavoro che gli
impedisce di continuare a lavorare. Lamenta un'importante invalidità e
vorrebbe essere risarcito dall'assicurazione. Siete consulenti per
l'assicurazione per valutare rischio di simulazione.

Il protagonista è un uomo di 26 anni che lavorava come magazziniere.
Mentre lavorava su un muletto, una grossa scatola è caduta e lo ha
colpito alla testa. Nell'impatto è caduto, ha perso il casco di
protezione e battuto la testa a terra, perdendo conoscenza. Portato al
pronto soccorso è stato confermato un trauma cranico e una serie di
fratture. Il paziente è stato sotto osservazione per alcuni giorni per
le fratture e poi dimesso.

Nei giorni successivi ha cominciato a lamentare disturbi di memoria e si
è rivolto alla sua assicurazione (che aveva stipulato 3 mesi prima) per
un risarcimento per potenziale invalidità.
\end{frame}

\begin{frame}{Risultati Valutazione Forense}
\phantomsection\label{risultati-valutazione-forense}
\tiny

\begin{longtable}[]{@{}
  >{\raggedright\arraybackslash}p{(\linewidth - 8\tabcolsep) * \real{0.3099}}
  >{\raggedright\arraybackslash}p{(\linewidth - 8\tabcolsep) * \real{0.2047}}
  >{\raggedright\arraybackslash}p{(\linewidth - 8\tabcolsep) * \real{0.1404}}
  >{\raggedright\arraybackslash}p{(\linewidth - 8\tabcolsep) * \real{0.1345}}
  >{\raggedright\arraybackslash}p{(\linewidth - 8\tabcolsep) * \real{0.2105}}@{}}
\toprule\noalign{}
\begin{minipage}[b]{\linewidth}\raggedright
Test Neuropsicologico
\end{minipage} & \begin{minipage}[b]{\linewidth}\raggedright
Dominio Valutato
\end{minipage} & \begin{minipage}[b]{\linewidth}\raggedright
Punteggio
\end{minipage} & \begin{minipage}[b]{\linewidth}\raggedright
Cut-off
\end{minipage} & \begin{minipage}[b]{\linewidth}\raggedright
Interpretazione
\end{minipage} \\
\midrule\noalign{}
\endhead
Trail Making Test -- Parte B & Flessibilità cognitiva & 240 sec & 180
sec & deficitario \\
WCST -- Errori perseverativi & Problem solving / astrazione & 32 & 20 &
deficitario \\
Torre di Londra -- Totale & Pianificazione & 5 & 18 & Deficitario \\
RAVLT -- Rich. immed. & Memoria verbale & 47 & 38 & Nella norma \\
Boston Naming Test (BNT) & Linguaggio & 49/60 & 48/60 & deficitario \\
SIMS -- Totale & Validità del profilo & 15 & 16 & No rischio
simulazione \\
\bottomrule\noalign{}
\end{longtable}
\end{frame}

\begin{frame}{}
\phantomsection\label{section}
\emph{Quali altre informazioni raccogliereste?}
\end{frame}

\begin{frame}{Approfondimenti vari nell'esempio}
\phantomsection\label{approfondimenti-vari-nellesempio}
Il paziente da quando ha smesso di lavorare ha cominciato a passare gran
parte del tempo a giocare a videogiochi e a giocare d'azzardo (online)
in camera. Soprattutto sembra passare gran parte del tempo giocando a
poker, vincendo e perdendo a seconda delle giornate.

Il paziente ha punteggi molto bassi anche quando corretti per età e
scolarità.
\end{frame}

\section{Conclusioni}\label{conclusioni}

\begin{frame}{Il ruolo della/del neuropsicologa/o nella valutazione}
\phantomsection\label{il-ruolo-delladel-neuropsicologao-nella-valutazione}
Vista la rilevanza delle interpretazioni nell'utilizzo dei test la/il
neuropsicologa/o ha un ruolo chiave.

\begin{itemize}
\item
  Deve conoscere le proprietà dei test che utilizza per interpretarli in
  maniera sensata
\item
  Deve avere un solido background in neuropsicologia clinica,
  neuroanatomia, per poter raccogliere informazioni durante la
  valutazione (in filosofia, ``teoreticità dell'osservazione'')
\end{itemize}
\end{frame}

\begin{frame}{Valutazione e ruolo attivo}
\phantomsection\label{valutazione-e-ruolo-attivo}
\pandocbounded{\includegraphics[keepaspectratio]{Figures/Active_collection_2.png}}
\end{frame}




\end{document}
